\documentclass[12pt,a4paper]{article}
\usepackage[T2A]{fontenc}
\usepackage[utf8]{inputenc}
\usepackage[russian]{babel}
\usepackage[left=2cm,right=2cm,top=2cm,bottom=2.5cm]{geometry}
\usepackage{hyperref}
\usepackage{microtype}

\begin{document}
\hypersetup{pageanchor=false}
\begin{titlepage}
    \centering
    Федеральное государственное автономное \\
    образовательное учреждение высшего образования \\
    \textbf{<<Национальный исследовательский университет ИТМО>>} \\
    \vfill
    Лабораторная работа \\
    \vspace{0.5\baselineskip}
    \textbf{Получение и печать Backtrace в xv6} \\
    \textbf{<<Операционные системы>>} \\
    \vspace{0.5\baselineskip}
    Вариант: Backtrace \\
    \vfill
    \flushright
    Выполнил:\\
    Лесников Владимир Алексеевич (Р3314) \\
    \vspace{0.5\baselineskip}
    Проверил:\\
    Смирнов Виктор Игоревич \\
    \vspace{3\baselineskip}
    \centering
    г. Санкт-Петербург, 2026
\end{titlepage}
\hypersetup{pageanchor=true}
\tableofcontents
\newpage

\section{Цель и задание}

Реализовать в ядре xv6 получение цепочки вызовов (backtrace) для стека ядра и пользовательского стека. Результат должен представлять собой массив адресов. Необходимо продемонстрировать печать backtrace в терминал, добавить её при вызове \texttt{panic} и написать скрипт преобразования адресов в имена функций.

Работа выполнена в учебной ОС xv6 для RISC-V. Запуск и проверка проводились в QEMU. Отдельный тип для backtrace в условии рекомендован, но не обязателен: в реализации используются массивы \texttt{uint64} и функции их заполнения и печати.

\section{Принцип получения backtrace}

При вызове функции её стековый кадр может содержать указатель на предыдущий кадр и адрес возврата. Последовательно читая эти записи, можно восстановить цепочку вызовов от более поздних к более ранним.

В используемом соглашении о вызовах RISC-V указатель кадра хранится в регистре \texttt{s0}. Предыдущий указатель находится по смещению \texttt{-16}, адрес возврата --- по смещению \texttt{-8} от указателя кадра. В проекте уже включён флаг компиляции \texttt{-fno-omit-frame-pointer}, сохраняющий кадровые указатели.

Ядро напрямую читает собственный стек. Пользовательский стек находится в другом адресном пространстве: исходные значения регистров берутся из \texttt{trapframe}, а чтение памяти выполняется после преобразования адреса через таблицу страниц процесса.

\section{Реализация}

В файле \path{kernel/printk.c} добавлены три процедуры:
\begin{itemize}
    \item \texttt{kernel\_backtrace} заполняет массив адресами возврата из кадров стека ядра и возвращает количество записей;
    \item \texttt{user\_backtrace} получает адреса из пользовательского стека текущего процесса;
    \item \texttt{print\_backtrace} печатает метку цепочки и каждый адрес отдельной строкой.
\end{itemize}

Обход ограничен ёмкостью массива. Проверяются выравнивание кадров и продвижение к старшим адресам. Пользовательская память проверяется через \texttt{walkaddr}; отсутствующие страницы при диагностике не выделяются. Первый пользовательский адрес берётся из сохранённого \texttt{ra}, поскольку ассемблерная обёртка системного вызова не создаёт обычный кадр.

В \texttt{panic} добавлен сбор и вывод backtrace ядра перед окончательной остановкой. Для демонстрации добавлен системный вызов \texttt{backtrace()}: обработчик в \path{kernel/sysproc.c} собирает и печатает обе цепочки, используя массив из 32 адресов. Номер и обработчик зарегистрированы в \path{kernel/syscall.h} и \path{kernel/syscall.c}; прототипы добавлены в \path{kernel/defs.h} и \path{user/user.h}, обёртка --- в \path{user/usys.pl}.

Программа \path{user/backtracetest.c} создаёт цепочку
\begin{center}
\texttt{main} $\rightarrow$ \texttt{trace\_middle} $\rightarrow$ \texttt{trace\_inner} $\rightarrow$ \texttt{backtrace}.
\end{center}
Для вспомогательных функций задан атрибут \texttt{noinline}, чтобы компилятор не встроил их в вызывающий код. Программа включена в образ файловой системы через список \texttt{UPROGS} в \path{Makefile}.

Скрипт \path{symbolize.sh} принимает ELF-файл и адреса в аргументах либо через стандартный ввод. Он вызывает \texttt{riscv64-elf-addr2line}, который по отладочной информации определяет имена функций и строки исходного кода. Другой инструмент можно задать переменной \texttt{ADDR2LINE}.

\section{Проверка работы}

Ядро и демонстрационная программа успешно собраны командой:
\begin{verbatim}
make kernel/kernel user/_backtracetest
\end{verbatim}

После запуска \texttt{make qemu} в оболочке xv6 выполнена команда \texttt{backtracetest}. Получен следующий вывод:
\begin{verbatim}
kernel backtrace:
0x0000000080002b08
0x0000000080002ac6
0x0000000080002850
0x0000003ffffff09c
user backtrace:
0x0000000000000010
0x0000000000000028
0x0000000000000040
0x0000000000000060
backtrace result: 2
\end{verbatim}

Символизация пользовательских адресов выполнена на хосте:
\begin{verbatim}
sh symbolize.sh user/_backtracetest 0x10 0x28 0x40 0x60
\end{verbatim}
Получены функции \texttt{trace\_inner}, \texttt{trace\_middle}, \texttt{main} и \texttt{start}. Также проверен ввод адресов через стандартный поток:
\begin{verbatim}
printf '%s\n' 0x10 0x28 0x40 0x60 |
  sh symbolize.sh user/_backtracetest
\end{verbatim}

Для проверки аварийного вывода внутри тестовой xv6 выполнена команда \texttt{kill 1}. Завершение процесса \texttt{init} запрещено, поэтому ядро вызвало \texttt{panic}:
\begin{verbatim}
panic: init exiting
kernel backtrace:
0x0000000080000a86
0x0000000080002266
0x00000000800028a6
0x0000003ffffff09c
\end{verbatim}

Первые три адреса расшифрованы командой:
\begin{verbatim}
sh symbolize.sh kernel/kernel 0x80000a86 0x80002266 0x800028a6
\end{verbatim}
Результат: \texttt{panic}, \texttt{kexit}, \texttt{usertrap}. Таким образом, печать при реальном вызове \texttt{panic} подтверждена. Выход из QEMU выполняется сочетанием \texttt{Ctrl+A}, затем \texttt{X}. После пересборки численные значения адресов могут измениться.

\section{Ограничения и результат}

Выполнены все обязательные пункты задания: получение обоих backtrace в виде массивов, отдельная процедура печати, демонстрация её работы, вывод при \texttt{panic} и скрипт символизации.

Обход ограничен ёмкостью массива. При демонстрации сохраняется до 32 адресов. Оптимизации и функции без кадров могут сокращать цепочку. Повреждённые стеки и одновременные вызовы \texttt{panic} на нескольких процессорах отдельно не проверялись.

\end{document}